\documentclass[10pt,a4paper]{amsart}
\usepackage[margin=19mm]{geometry}
\usepackage{amsmath,amssymb,mathtools}
\usepackage[scr=rsfs,cal=euler]{mathalfa}
\usepackage{xfrac}
\usepackage[unicode,hidelinks,bookmarks=false]{hyperref}
\newcommand{\ie}{i.e.\ }
\newcommand{\cf}{cf.\ }
\newcommand{\resp}{respectively\ }
\newcommand{\Subsection}{\subsection}
\providecommand{\nobreakdash}{\nobreak\hbox{-}}
\setlength{\emergencystretch}{2em}
\multlinegap0pt
\usepackage{bm}
\usepackage{bbm}
\usepackage{dsfont}
\usepackage{xspace}
\usepackage{cancel}
\usepackage{mathtools}
\usepackage{amsthm}
\makeatletter
\@ifundefined{theorem}{\newtheorem{theorem}{Theorem}}{}
\@ifundefined{lemma}{\newtheorem{lemma}[theorem]{Lemma}}{}
\makeatother
\DeclareMathOperator{\dist}{dist}
\newcommand{\Om}{\Omega}
\newcommand{\R}{\mathbb R}
\newcommand{\eps}{\varepsilon}
\newcommand{\mm}{\mathfrak m}
\title[Arbitrarily Many Non-degenerate Local Maxima of First Nonzer]{Arbitrarily Many Non-degenerate Local Maxima of First Nonzero Eigenfunctions on Positively Curved Two-spheres}
\author{Heng Zhang}
\date{Source: arXiv:2606.26419v1 (CC BY 4.0)}
\begin{document}
\maketitle

\noindent\emph{Source and licence.} Arbitrarily Many Non-degenerate Local Maxima of First Nonzero Eigenfunctions on Positively Curved Two-spheres. Version arXiv:2606.26419v1 (CC BY 4.0), distributed under the Creative Commons Attribution 4.0 International licence, as recorded in the arXiv licence metadata for this exact version and shown on the abstract page https://arxiv.org/abs/2606.26419v1. Full rights notice: licenses/arxiv-2606.26419-CC-BY-4.0.txt.\par\smallskip

\noindent\textit{Editorial scope.} Counted unit: the Lemma whose source label is \texttt{lem:pancake-convergence-fixed-delta} in the article (source0.tex lines 1257--1265, source label \texttt{lem:pancake-convergence-fixed-delta}), delivered together with the article's own complete original proof of it (source0.tex lines 1267--1378, one \texttt{proof} environment of 112 lines). No mathematics is written, repaired or re-derived: statement and proof are the article's own text line for line. Required context retained uncounted: lem:positive-curvature (lines 1212--1216). Every definition, display and label of those lines is retained unchanged. Omitted from this package and not redistributed: 1--1211, 1217--1256, 1266--1266, 1379--2078. Nothing inside a retained excerpt is omitted or altered: every retained body is delivered exactly as the article's lines are written, so no omission receipt of any kind accompanies this package. Citations: the retained excerpts carry no citation command at all, so no bibliography entry of the article is transcribed and no reference is redistributed. Reference adaptation: none was needed and none was made; every reference inside the delivered unit resolves inside it, so no unresolved reference or citation is produced. Printed slips kept literal: no printed word, symbol, hypothesis or number of the counted statement or of its proof was altered. Layout and class: the article's own class is \texttt{amsart} with packages including geometry, amsmath, amssymb, amsfonts, amsthm, mathtools, mathrsfs, enumitem, hyperref, tikz; the wrapper is the lane's pinned \texttt{amsart} editorial wrapper at 10pt on A4, which restates the theorem-like environments the retained text needs and adds the article's own macro definitions that the retained text needs (\texttt{\textbackslash Om}, \texttt{\textbackslash R}, \texttt{\textbackslash dist}, \texttt{\textbackslash eps}, \texttt{\textbackslash mm}), with extra wrapper packages (bm, bbm, dsfont, xspace, cancel, mathtools). The delivered numbering is the unit's own. Assets: no retained span contains a figure environment, a graphics-inclusion command, a PDF-page inclusion, a table or a TikZ picture; no image file is copied into this package, the package declares no asset dependency and no image file is redistributed with it. Licence: version arXiv:2606.26419v1 of this article is distributed under the Creative Commons Attribution 4.0 International licence, as recorded in the arXiv licence metadata for this exact version and shown on the abstract page https://arxiv.org/abs/2606.26419v1. A full rights notice is delivered at licenses/arxiv-2606.26419-CC-BY-4.0.txt. Characters: every retained excerpt is pure ASCII, so the delivered engine, XeLaTeX with its default Latin Modern text font, reports no missing character.
\begin{lemma}\label{lem:positive-curvature}
For every sufficiently small $\delta>0$ and every $\eps>0$,
$M_{\eps,\delta}$ is a smooth embedded two-sphere with positive Gaussian
curvature.
\end{lemma}

\begin{lemma}\label{lem:pancake-convergence-fixed-delta}
For each fixed sufficiently small $\delta>0$,
\begin{equation}\label{eq:pancake-fixed-delta-convergence}
(M_{\eps,\delta},d_{M_{\eps,\delta}},\mm_{M_{\eps,\delta}})
\longrightarrow
(D\Om_\delta,d_{D\Om_\delta},\mm_{D\Om_\delta})
\end{equation}
in measured Gromov--Hausdorff sense as $\eps\downarrow0$.
\end{lemma}

\begin{proof}
Set
$$
f_\delta(x)=\sqrt{1-\psi_\delta(x)},
\qquad x\in\overline{\Om_\delta}.
$$
Because $\psi_\delta$ is convex, $1-\psi_\delta$ is concave and nonnegative on
$\Om_\delta$; since $t\mapsto\sqrt t$ is increasing and concave on
$[0,\infty)$, $f_\delta$ is concave.  Also $0\le f_\delta\le1$.
Note that $M_{\eps,\delta}$ is the union of the two graphs
\begin{equation}\label{eq:pancake-two-graphs}
z=\eps f_\delta(x),
\qquad
z=-\eps f_\delta(x),
\qquad
x\in\overline{\Om_\delta}.
\end{equation}
The graph parametrization in \eqref{eq:pancake-two-graphs} is singular at
$\partial\Om_\delta$, but Lemma~\ref{lem:positive-curvature} proves that the
level surface itself is smooth there.

Define
$$
\pi_{\eps,\delta}:M_{\eps,\delta}\to D\Om_\delta
$$
by horizontal projection, sending the upper graph to the positive sheet and the
lower graph to the negative sheet.  This is well-defined on the seam because
$f_\delta=0$ on $\partial\Om_\delta$, and the two boundary copies are identified
in the double.  Horizontal projection does not increase the length of rectifiable
curves.  Therefore
\begin{equation}\label{eq:pancake-projection-lower-bound}
d_{D\Om_\delta}(\pi_{\eps,\delta}(p),\pi_{\eps,\delta}(q))
\le d_{M_{\eps,\delta}}(p,q)
\qquad\text{for all }p,q\in M_{\eps,\delta}.
\end{equation}

Conversely, suppose first that $p,q$ lie on the same sheet and that
$\pi_{\eps,\delta}(p)=x^\sigma$, $\pi_{\eps,\delta}(q)=y^\sigma$.  The Euclidean
segment $\gamma(t)=(1-t)x+ty$ is contained in $\Om_\delta$.  Since
$f_\delta\circ\gamma$ is concave, it has bounded variation, and the length of its
lift to the corresponding graph satisfies
\begin{equation}\label{eq:same-sheet-lift-length}
\mathcal L\bigl(t\mapsto(\gamma(t),\pm\eps f_\delta(\gamma(t)))\bigr)
\le |x-y|+\eps\operatorname{Var}(f_\delta\circ\gamma).
\end{equation}
Because $0\le f_\delta\le1$, the total variation of the one-dimensional concave
function $f_\delta\circ\gamma$ is at most $2$.  Thus the lifted segment in
\eqref{eq:same-sheet-lift-length} has length at most $|x-y|+2\eps$.

If $p$ and $q$ lie on opposite sheets, then, for arbitrary $\rho>0$, choose
$a\in\partial\Om_\delta$ such that
$$
|x-a|+|a-y|\le d_{D\Om_\delta}(x^+,y^-)+\rho.
$$
Lift $x\to a$ on the upper graph and $a\to y$ on the lower graph.  By the same
estimate, the total lifted length is at most
$ d_{D\Om_\delta}(x^+,y^-)+4\eps+\rho $.  Letting $\rho\downarrow0$ gives
\begin{equation}\label{eq:pancake-projection-upper-bound}
d_{M_{\eps,\delta}}(p,q)
\le d_{D\Om_\delta}(\pi_{\eps,\delta}(p),\pi_{\eps,\delta}(q))+4\eps.
\end{equation}
Equations \eqref{eq:pancake-projection-lower-bound} and
\eqref{eq:pancake-projection-upper-bound} show that $\pi_{\eps,\delta}$ is onto
and has distortion at most $4\eps$.  This is Gromov--Hausdorff convergence.

It remains to prove measure convergence.  On one graph, away from the boundary,
the area density is
\begin{equation}\label{eq:area-density}
dA_{\eps,\delta}=\sqrt{1+\eps^2|\nabla f_\delta|^2}\,dx.
\end{equation}
Although $f_\delta$ is not Lipschitz up to $\partial\Om_\delta$, the formula
\eqref{eq:area-density} is justified by applying the graph area formula first on
the compact exhaustion
$$
\Om_{\delta,s}:=\{x\in\Om_\delta:\dist_{\R^2}(x,\partial\Om_\delta)\ge s\},
$$
where $f_\delta$ is smooth, and then letting $s\downarrow0$.  The exhausted graph
pieces increase to the whole graph up to the seam, the seam itself has zero
two-dimensional area, and monotone convergence of the areas applies once the
finite collar contribution is verified by the estimate below.

We claim that
\begin{equation}\label{eq:grad-f-integrable}
|\nabla f_\delta|\in L^1(\Om_\delta).
\end{equation}
Away from the boundary this is clear.  Near $\partial\Om_\delta$, the boundary is
smooth and $|\nabla\psi_\delta|$ is bounded above and below by positive
constants.  In a normal collar, if $r=\dist_{\R^2}(x,\partial\Om_\delta)$, Taylor
expansion gives $1-\psi_\delta(x)\simeq c r$ and
$|\nabla\psi_\delta(x)|\le C$.  Hence
$$
|\nabla f_\delta(x)|=
\frac{|\nabla\psi_\delta(x)|}{2\sqrt{1-\psi_\delta(x)}}
\le C r^{-1/2},
$$
which is integrable in a two-dimensional collar.  This proves
\eqref{eq:grad-f-integrable}.  Therefore, for $0<\eps\le1$,
$$
0\le \sqrt{1+\eps^2|\nabla f_\delta|^2}-1\le \eps |\nabla f_\delta|
\le |\nabla f_\delta|,
$$
and dominated convergence, using \eqref{eq:grad-f-integrable}, gives
$$
\sqrt{1+\eps^2|\nabla f_\delta|^2}\to1
\quad\text{in }L^1(\Om_\delta).
$$
The same holds on the lower graph.  Thus
$(\pi_{\eps,\delta})_\#\mm_{M_{\eps,\delta}}$ converges weakly to the sum
of two copies of Lebesgue measure on $\Om_\delta$, namely to
$\mm_{D\Om_\delta}$.  Together with the Gromov--Hausdorff convergence above,
this proves \eqref{eq:pancake-fixed-delta-convergence}.
\end{proof}

\end{document}
